\chapter{Streaming arithmetic with continued fractions}\label{app:gosper}
\chaptertag{Gosper's homographic and bihomographic algorithms}

Suppose a number is arriving one continued-fraction digit at a time.  Can we
transform it without first converting it to a decimal, and without waiting for
the entire infinite input?  Bill Gosper's answer is yes.  The method reads just
enough of the input to prove the next output digit, emits that digit, and then
continues.

\section{The problem---before the machinery}

There are two versions.

\begin{description}[leftmargin=2.7cm,style=nextline]
  \item[One input.] Given
  \[
    x=[a_0;a_1,a_2,\dots]
  \]
  as a stream, and a linear fractional transformation
  \[
    L(x)=\frac{\alpha x+\beta}{\gamma x+\delta},
  \]
  produce the continued fraction of $L(x)$ from left to right.

  \item[Two inputs.] Given continued-fraction streams for $x$ and $y$, produce
  the continued fraction of a two-variable fractional expression.  In
  particular, $x+y$, $x-y$, $xy$, and $x/y$ all have the form
  \[
    \Phi(x,y)=\frac{axy+bx+cy+d}{exy+fx+gy+h}.
  \]
\end{description}

The intended computation is \emph{exact}.  It should never guess a digit from
a floating-point approximation.  Instead, every emitted digit comes with a
proof that all inputs compatible with the digits seen so far would produce that
same digit.

\begin{bigidea}
A finite input prefix determines an \emph{interval} of possible values.  Map
that interval through the desired transformation.  If the whole image fits
inside one continued-fraction digit cell $[k,k+1)$, the next output digit is
certified to be $k$.  If it crosses an integer boundary, read more input and
shrink the interval.
\end{bigidea}

The original algorithm appeared as Item 101B of the 1972 MIT memo
\emph{HAKMEM}.  A modern refinement by Michael J. Collins allows exact
intervals themselves to appear in the output stream, resolving boundary cases
in which no next digit can be certified after finite time.

\section{What the unread tail means}

The subscript on a tail records the \emph{next} digit, not the first digit of
the original number.  After the prefix $a_0,\dots,a_{j-1}$ has been read, write
\begin{equation}
  x=[a_0;a_1,\dots,a_{j-1},t_j],
  \qquad
  t_j=[a_j;a_{j+1},\dots].                         \label{eq:tail-invariant}
\end{equation}
Here $t_j$ is the complete unread remainder of the continued fraction, used as
the last entry in the first expression.  Before $a_j$ is known, the only
universal information is
\[
  1\leq t_j\leq\infty
\]
for a positive simple continued fraction.  When the next digit $a_j=m$ arrives,
\begin{equation}
  t_j=m+\frac1{t_{j+1}},
  \qquad t_{j+1}=[a_{j+1};a_{j+2},\dots]\in[1,\infty].
  \label{eq:tail-shift}
\end{equation}
Thus knowing the single digit $m$ restricts the old tail to
$t_j\in(m,m+1]$; absorbing $m$ replaces it by the new unknown tail
$t_{j+1}\in[1,\infty]$.

\begin{labbox}[title={Known prefix versus unknown tail}]
At the start, $t_0=x$ and no input digits have been absorbed.  After $a_0$ is
read, the state incorporates $a_0$ and the unknown variable is
$t_1=[a_1;a_2,\dots]$.  After $a_1$ is read, the state incorporates
$a_0,a_1$ and the unknown variable is $t_2$, and so on.  The state changes;
the exact number represented does not.
\end{labbox}

Let
\[
  D(m)=\begin{pmatrix}m&1\\1&0\end{pmatrix},
  \qquad
  M_0=\begin{pmatrix}\alpha&\beta\\\gamma&\delta\end{pmatrix}.
\]
After $j$ input digits have been absorbed, set
\begin{equation}
  M_j=M_0D(a_0)D(a_1)\cdots D(a_{j-1})
      =\begin{pmatrix}A&B\\C&D\end{pmatrix}.       \label{eq:prefix-state}
\end{equation}
Then the output value still waiting to be written is exactly
\[
  F_{M_j}(t_j)=\frac{At_j+B}{Ct_j+D}.
\]
This is the invariant behind the algorithm.

\section{The geometry: map an interval, then inspect its floor}

For a general certified tail interval $I=[b_0,b_1]$, consider
\[
  F_M(t)=\frac{At+B}{Ct+D}.
\]
Provided $Ct+D$ never vanishes on $I$, the derivative
\begin{equation}
  F_M'(t)=\frac{AD-BC}{(Ct+D)^2}                    \label{eq:mobius-derivative-b}
\end{equation}
has constant sign.  Hence $F_M$ is monotone, and the image $F_M(I)$ is simply
the interval whose endpoints are $F_M(b_0)$ and $F_M(b_1)$, in the appropriate
order.  No point between them can escape those endpoint bounds.

\begin{center}
\begin{tikzpicture}[x=1cm,y=1cm]
  \draw[navy!55,line width=.8pt] (0,1.25)--(5.2,1.25);
  \fill[teal] (1,1.25) circle (2.5pt) node[below=5pt,text=ink] {$b_0$};
  \fill[teal] (4.2,1.25) circle (2.5pt) node[below=5pt,text=ink] {$b_1$};
  \draw[teal,line width=3pt] (1,1.25)--(4.2,1.25);
  \node[text=muted] at (2.6,1.75) {possible input tail $I$};

  \draw[-{Stealth[length=2.5mm]},coral,line width=1.2pt]
    (5.65,1.25) to[bend left=22] (7.1,1.25);

  \foreach \x/\lab in {7.4/$k$,9.0/$k+1$,10.6/$k+2$}{
    \draw[navy!35] (\x,.94)--(\x,1.58);
    \node[below=4pt,text=muted] at (\x,.94) {\lab};
  }
  \draw[gold,line width=5pt] (7.75,1.25)--(8.65,1.25);
  \fill[coral] (7.75,1.25) circle (2.4pt);
  \fill[coral] (8.65,1.25) circle (2.4pt);
  \node[text=muted] at (8.2,1.75) {$F_M(I)\subset[k,k+1)$};
\end{tikzpicture}

\small\color{muted}If both endpoint images lie in the same floor cell and
there is no pole between them, every possible output has floor $k$.
\end{center}

For the universal interval $[1,\infty]$, the two endpoint values are
\[
  F_M(1)=\frac{A+B}{C+D},
  \qquad
  F_M(\infty)=\frac AC.
\]
An endpoint at infinity may be open, and a finite continued fraction may end;
exact implementations retain this endpoint information.  The essential test is
unchanged:
\[
  F_M(I)\subseteq[k,k+1)quad\Longrightarrow\quad
  \text{the next output digit is }k.
\]
If the image crosses an integer boundary, or if a pole is possible, no digit is
yet justified.  We must refine the input interval.

\section{The two algebraic moves}

\paragraph{Absorb the next input digit.}
If $a_j=m$, substitute $t_j=m+1/t_{j+1}$ into the current map.  Clearing the
factor $t_{j+1}$ gives
\[
 F_M\left(m+\frac1{t_{j+1}}\right)
 =\frac{(Am+B)t_{j+1}+A}{(Cm+D)t_{j+1}+C}.
\]
Thus the exact state update is
\begin{equation}
  \boxed{\mathcal I_m(M)=M D(m)
  =\begin{pmatrix}Am+B&A\\Cm+D&C\end{pmatrix}.}    \label{eq:unary-absorb}
\end{equation}

\paragraph{Emit a certified output digit.}
If the image interval lies in $[k,k+1)$, write the digit $k$.  The unread part
of the output continued fraction is then $1/(F_M-k)$, represented by
\begin{equation}
  \boxed{\mathcal E_k(M)=
  \begin{pmatrix}0&1\\1&-k\end{pmatrix}M
  =\begin{pmatrix}C&D\\A-kC&B-kD\end{pmatrix}.}   \label{eq:unary-emit}
\end{equation}

\begin{checkpoint}[title={The homographic loop}]
Maintain the statement: ``the unread output is $F_M(t_j)$ for every
$t_j\in I$.''
\begin{enumerate}
  \item Map the current certified interval $I$ through $F_M$ by evaluating its
  endpoints.
  \item If the image lies in one floor cell $[k,k+1)$, emit $k$, apply
  $M\leftarrow\mathcal E_k(M)$, and replace the output interval by
  $1/(F_M(I)-k)$.
  \item Otherwise read $a_j=m$, use $t_j=m+1/t_{j+1}$, apply
  $M\leftarrow\mathcal I_m(M)$, and return to the universal interval
  $t_{j+1}\in[1,\infty]$.
  \item Repeat.  A common nonzero factor may be removed from the four matrix
  entries without changing the map.
\end{enumerate}
\end{checkpoint}

\section{Worked stream: the first digits of \texorpdfstring{$\pi/2$}{pi/2}}

Here $L(x)=x/2$, so
$M_0=\begin{psmallmatrix}1&0\\0&2\end{psmallmatrix}$, and the input is
\[
 \pi=[3;7,15,1,292,\dots].
\]
After reading $a_0=3$ and $a_1=7$, the unknown variable is
$t_2=[15;1,292,\dots]\in[1,\infty]$, and
\[
 M_2=M_0D(3)D(7)=\begin{pmatrix}22&3\\14&2\end{pmatrix}.
\]
The first row below maps the endpoint tails $1$ and $\infty$.  Each following
row applies \eqref{eq:unary-emit}; no further input digit is read.

\begin{center}
\small
\begin{tabular}{c c c}
\toprule
event & current matrix $M$ & image of the possible tail interval \\
\midrule
read $3,7$ & $\begin{psmallmatrix}22&3\\14&2\end{psmallmatrix}$
 & $[25/16,\,11/7)$: emit $1$ \\
emit $1$ & $\begin{psmallmatrix}14&2\\8&1\end{psmallmatrix}$
 & $(7/4,\,16/9]$: emit $1$ \\
emit $1$ & $\begin{psmallmatrix}8&1\\6&1\end{psmallmatrix}$
 & $[9/7,\,4/3)$: emit $1$ \\
emit $1$ & $\begin{psmallmatrix}6&1\\2&0\end{psmallmatrix}$
 & $(3,\,7/2]$: emit $3$ \\
\bottomrule
\end{tabular}
\end{center}

After these four emissions the state is
$\begin{psmallmatrix}2&0\\0&1\end{psmallmatrix}$: the unread output is
$2t_2$.  Only now do we read the next input digits, $15$ and $1$; absorbing
them maps the new unknown tail into $[31,32)$ and certifies $31$.  Therefore
\[
 \frac\pi2=[1;1,1,3,31,\dots].
\]
The calculation has used exact rational endpoints throughout.  The decimal
digits of $\pi$ never entered.

\section{Two inputs: rectangles replace intervals}

Now suppose prefixes of two inputs have been read.  Their unread tails, say
$u_j$ and $v_\ell$, range independently over certified intervals $I$ and $J$.
The unread output has the form
\begin{equation}
 B_{\mathcal B}(u,v)
 =\frac{auv+bu+cv+d}{euv+fu+gv+h},
 \qquad
 \mathcal B=(a,b,c,d;\,e,f,g,h).                  \label{eq:bihomographic}
\end{equation}
Thus the possible input is a rectangle $I\times J$, not a line segment.

\begin{center}
\begin{tikzpicture}[x=1cm,y=1cm]
  \fill[teal!10] (.6,.35) rectangle (4.2,2.45);
  \draw[teal,line width=1pt] (.6,.35) rectangle (4.2,2.45);
  \foreach \p/\lab/\pos in {
    {.6,.35}/{$(b_0,c_0)$}/{below left},
    {4.2,.35}/{$(b_1,c_0)$}/{below right},
    {.6,2.45}/{$(b_0,c_1)$}/{above left},
    {4.2,2.45}/{$(b_1,c_1)$}/{above right}}
    \fill[coral] (\p) circle (2.5pt) node[\pos,text=ink] {\lab};
  \node[text=muted] at (2.4,1.4) {$I\times J$};
  \draw[-{Stealth[length=2.5mm]},coral,line width=1.2pt]
    (4.85,1.4) to[bend left=18] (6.35,1.4);
  \draw[navy!55] (6.65,1.4)--(10.8,1.4);
  \draw[gold,line width=5pt] (7.35,1.4)--(9.85,1.4);
  \fill[coral] (7.35,1.4) circle (2.5pt);
  \fill[coral] (9.85,1.4) circle (2.5pt);
  \node[text=muted] at (8.6,1.95) {range bounded by corner images};
\end{tikzpicture}
\end{center}

Assume the denominator has no zero on the rectangle.  For a fixed $v$, the
map is linear fractional in $u$, so its extrema along that horizontal segment
occur at the two ends.  On either vertical edge it is linear fractional in
$v$, so the remaining extrema occur at the edge endpoints.  Consequently the
entire range is bounded by the images of the four corners.  In particular, for
$I=J=[1,\infty]$ these values are
\begin{equation}
 \frac ae,\qquad \frac{a+b}{e+f},\qquad
 \frac{a+c}{e+g},\qquad
 \frac{a+b+c+d}{e+f+g+h},                         \label{eq:four-corners}
\end{equation}
with the usual interpretation of infinite and open endpoints.

The four basic operations differ only in their starting coefficients:
\[
\begin{array}{c|c}
\text{desired output} & (a,b,c,d;\,e,f,g,h)\\ \hline
x+y &(0,1,1,0;\,0,0,0,1)\\
x-y &(0,1,-1,0;\,0,0,0,1)\\
xy  &(1,0,0,0;\,0,0,0,1)\\
x/y &(0,1,0,0;\,0,0,1,0)
\end{array}
\]

If the next $u$-digit is $m$, substitute $u=m+1/u'$; if the next $v$-digit is
$n$, substitute $v=n+1/v'$.  Collecting coefficients gives
\begin{align}
\mathcal I^u_m(\mathcal B)
 &=(ma+c,\,mb+d,\,a,\,b;\,me+g,\,mf+h,\,e,\,f),
                                                        \label{eq:binary-absorb-u}\\
\mathcal I^v_n(\mathcal B)
 &=(na+b,\,a,\,nc+d,\,c;\,ne+f,\,e,\,ng+h,\,g).
                                                        \label{eq:binary-absorb-v}
\end{align}
When the four corner images all lie in one floor cell $[k,k+1)$, emit $k$ and
invert the fractional part:
\begin{equation}
\mathcal E_k(\mathcal B)
=(e,f,g,h;\,a-ke,b-kf,c-kg,d-kh).                  \label{eq:binary-emit}
\end{equation}

The geometry and the logic are exactly as before.  If the image of the current
rectangle crosses an integer boundary, absorb from one input and replace the
rectangle by a thinner one; alternate inputs, or choose the one expected to
shrink the image more.  If the denominator may vanish inside the rectangle,
no output is safe until further input or separate algebra resolves the issue.

\begin{bigidea}
The four coefficients in the one-input problem, or eight in the two-input
problem, are bookkeeping for a moving interval image.  Matrix formulas make
the updates fast, but the proof of correctness is geometric: every possible
tail lies in the current interval or rectangle, and every emitted digit is
shared by its entire image.
\end{bigidea}

\section{The boundary problem and the interval repair}

The floor function jumps at every integer.  This creates a subtle failure mode.
Feed two \emph{independent} copies of the continued fraction of $\sqrt2$ to the
multiplication algorithm.  The actual product is $2$, but every finite pair of
independent enclosures allows one copy to move a little below $\sqrt2$ and the
other a little above it.  The resulting product rectangle can straddle $2$.
Its image then meets both floor cells $[1,2)$ and $[2,3)$, so the rule cannot
certify even the first digit.

This is an information problem, not numerical instability.  The two stream
variables have forgotten that they denote the same algebraic number.  Retaining
the relation $x=y$ reduces the expression to $x^2$, and the separate fact
$x^2=2$ resolves the equality immediately.

Collins's interval-valued refinement distinguishes a continued fraction from
its computational stream.  A stream item may be either
\begin{itemize}
  \item an explicit partial quotient, or
  \item an exact rational interval bounding the current unread value.
\end{itemize}
If no digit is forced, the algorithm may output a tighter certified interval
instead of waiting silently.  A consumer asks for an error tolerance
$\varepsilon$ and reads until the resulting value interval has width below
$\varepsilon$.  Exact endpoints and their open-or-closed status preserve the
no-roundoff character of Gosper's method while guaranteeing useful progress at
an integer boundary.

\begin{warningbox}[title={Digits and numerical accuracy are different requests}]
There may be no finite proof of the next continued-fraction digit even when a
very accurate numerical enclosure is easy.  A robust implementation therefore
asks for a certified error tolerance, not an unconditional number of output
digits.  Away from boundaries, intervals quickly collapse into ordinary
partial quotients and recover Gosper's original stream.
\end{warningbox}

\section{Exercises}

\begin{exercisebox}
\begin{enumerate}
  \item Suppose only the prefix $[2;3]$ of $x$ is known.  Write
  $x=[2;3,t_2]$ with $t_2\in[1,\infty]$, evaluate the two endpoint values,
  and prove that $x\in[9/4,7/3)$. \exlevel{Fluency}

  \item Prove the endpoint-image rule using
  \eqref{eq:mobius-derivative-b}.  Apply it to
  $F(t)=(22t+3)/(14t+2)$ on $[1,\infty]$ and certify its floor.
  \exlevel{Reasoning}

  \item Derive the absorption and emission updates
  \eqref{eq:unary-absorb} and \eqref{eq:unary-emit} directly by
  substitution.  Express each as matrix multiplication. \exlevel{Fluency}

  \item Continue the $\pi/2$ calculation: from
  $\begin{psmallmatrix}2&0\\0&1\end{psmallmatrix}$, absorb $15$ and $1$ and
  show that the new image lies in $[31,32)$. \exlevel{Reasoning}

  \item Prove the four-corner range rule for a bihomographic map on a finite
  rectangle with no denominator zero.  Why is checking four corners not enough
  if a pole lies inside? \exlevel{Reasoning}

  \item Derive \eqref{eq:binary-absorb-u} and
  \eqref{eq:binary-absorb-v}.  Verify that the four initial states really
  represent $+$, $-$, $\times$, and $/$. \exlevel{Fluency}

  \item Explain why independent intervals for two copies of $\sqrt2$ can
  straddle the boundary $2$ forever when computing $\sqrt2\sqrt2$.  State the
  extra algebraic information that resolves the equality. \exlevel{Reasoning}

  \item Design a lazy implementation whose output type is either
  \texttt{Digit(k)} or \texttt{Bound(lo,hi)}.  State invariants that make every
  emitted item a certificate, then test all four operations on rational and
  quadratic inputs. \exlevel{Project}
\end{enumerate}
\end{exercisebox}
